\chapter{Athuls Contributions}

\section{Extended Literature Review}

\subsection{Early Detection and Temporal Dynamics}

A key limitation of content-based misinformation detection models is their reliance on fully propagated social context, which may not be available in early detection scenarios. Liu and Wu \cite{Liu2020} address this challenge through FNED, a deep network designed for fake news early detection on social media. FNED models the sequential behaviour of crowd responses over time, combining textual features with user-profile signals to classify news veracity within minutes of initial publication. Evaluated on Twitter and Weibo, FNED achieves over 90\% accuracy within the first five minutes of propagation, demonstrating that temporal behavioural patterns alone carry strong classification signals. This finding is particularly relevant to the MACE framework's Stage 2 experimental design, where early detection robustness under incomplete propagation data represents an important but underexplored evaluation dimension.

\subsection{Bi-Directional Propagation Modelling}

Expanding on graph-based approaches to social context modelling, Bian et al.\ \cite{Bian2020} propose Bi-Directional Graph Convolutional Networks (Bi-GCN) for rumour detection on social media. Unlike unidirectional propagation models, Bi-GCN simultaneously encodes both the top-down spreading pattern and the bottom-up aggregation of responses within a reply tree structure. Evaluated on Twitter15, Twitter16, and Weibo — the same benchmark datasets used in the MACE Stage 1 replication — Bi-GCN demonstrates that modelling bidirectional propagation significantly improves classification performance over prior graph-based baselines. This work provides additional motivation for the retention of the DANES and GETAE social context branches in the MACE framework, as these architectures similarly leverage structured propagation signals that single-direction models fail to capture fully.

\subsection{Image-Text Similarity and Cross-Modal Consistency}

Zhou et al.\ \cite{Zhou2020} introduce SAFE, a Similarity-Aware Multi-modal Fake News Detection framework that explicitly models the relationship between textual and visual content. SAFE extracts semantic representations from both text and images and computes a similarity score to identify inconsistencies between modalities — a common characteristic of manipulated or misleading posts. Evaluated on a Twitter-based dataset, SAFE outperforms prior unimodal and multimodal baselines, achieving strong accuracy while providing interpretable similarity scores. This approach is directly relevant to the multimodal content branch of MACE, as SAFE's cross-modal consistency modelling offers a complementary perspective to HMCAN's hierarchical attention fusion. Integrating similarity-based signals could further strengthen MACE's content branch, particularly for detecting posts where visual and textual content are semantically misaligned.

\subsection{Multimodal Variational Representation Learning}

Khattar et al.\ \cite{Khattar2019} propose MVAE, a Multimodal Variational Autoencoder for fake news detection that learns a shared latent representation from both textual and visual modalities through a bimodal variational autoencoder coupled with a binary classifier. Unlike discriminative models that rely on supervised feature extraction, MVAE employs generative modelling to capture the joint distribution of multimodal content, enabling the detection of subtle cross-modal inconsistencies that may not surface through attention-based fusion alone. Evaluated on Twitter and Weibo datasets, MVAE outperforms state-of-the-art methods by margins of approximately 6\% in accuracy and 5\% in F1 score. As a generative multimodal baseline, MVAE provides a useful point of comparison for the MACE framework's content branch evaluation, particularly in ablation settings where the contribution of visual features is isolated from social context signals.

\subsection{Out-of-Context Multimodal Misinformation}

Luo et al.\ \cite{Luo2021} present NewsCLIPpings, a large-scale benchmark dataset for out-of-context multimodal misinformation detection, constructed by automatically pairing real news images with semantically related but contextually mismatched captions. Unlike datasets where both text and image are fabricated, NewsCLIPpings specifically targets a prevalent real-world misinformation strategy in which authentic images are repurposed with misleading textual context. The dataset, drawn from four major news sources via the VisualNews corpus, has since become a standard evaluation benchmark for multimodal misinformation detection systems. Its inclusion as a supplementary evaluation dataset for MACE would strengthen the generalisability of findings beyond the Twitter-based benchmarks used in DANES and GETAE, and provide additional insight into the performance of the multimodal content branch when visual-textual consistency is deliberately compromised.

\subsection{Comprehensive Survey of Deep Learning Approaches}

Hu et al.\ \cite{Hu2022} provide a comprehensive survey of deep learning methods for fake news detection, systematically reviewing approaches spanning text-only content models, multimodal frameworks, social context architectures, and knowledge-enhanced systems. The survey identifies the integration of multimodal content with social propagation signals as an open research challenge, and highlights the absence of controlled ablation studies that isolate the contribution of each component as a key limitation of existing work. This observation directly motivates the comparative experimental design adopted by MACE, wherein controlled substitution of the content branch and ablation-based evaluation are employed to decompose the relative contributions of multimodal features and social context signals. The survey further confirms that the research questions addressed by this project — evaluating the impact of multimodal substitution within context-aware ensemble architectures — represent a genuine and unresolved gap in the current literature.

\subsection{Summary of Additional Literature Contributions}

The papers reviewed in this chapter complement the existing literature review by addressing dimensions not fully covered in the primary survey. Collectively, they reinforce three key arguments underlying the MACE framework: first, that temporal and propagation-based social context signals provide strong classification value independent of content modality \cite{Liu2020, Bian2020}; second, that cross-modal consistency and shared multimodal representations enhance fake news detection beyond what either modality achieves independently \cite{Zhou2020, Khattar2019}; and third, that the controlled evaluation of multimodal substitution within context-aware architectures represents a recognised and unaddressed gap in the literature \cite{Luo2021, Hu2022}. Together, these contributions strengthen the motivation for the MACE framework and provide additional benchmarking context for its experimental evaluation.
